\section{Application of Program Outcomes}
Table~\ref{tab:po} maps each Program Outcome and Program Specific Outcome addressed by this case study to its concrete implementation. Every mapping demonstrates how foundational engineering knowledge, modern tools, and professional ethics are exercised within the project.

{\small
\begin{longtable}{@{}>{\centering\arraybackslash}p{1.8cm} p{13.2cm}@{}}
    \caption{Mapping of Program Outcomes and Program Specific Outcomes to the Case Study.}\label{tab:po}\\
    \toprule
    \textbf{Outcome} & \textbf{Outcome Description and Concrete Relevance to the Case Study} \\
    \midrule
    \endfirsthead
    \toprule
    \textbf{Outcome} & \textbf{Outcome Description and Concrete Relevance to the Case Study} \\
    \midrule
    \endhead
\textbf{PO3} & \textbf{Design/Development of Solutions:} Designs scalable system architectures, robust database schemas, object hierarchies, or automated pipelines meeting specified operational needs with proper consideration for efficiency, maintainability, and safety. In this case study (Time-Series Forecasting Using LSTM Networks), a multi-stage deep learning pipeline spanning tensor augmentation, backbone feature extraction, and optimized inference heads is designed. \\
    \addlinespace
\textbf{PO1} & \textbf{Engineering Knowledge:} Applies principles of mathematics, computer science, and engineering fundamentals to formalize, model, and analyze the system architecture, mathematical constraints, and foundational algorithms. In this case study (Time-Series Forecasting Using LSTM Networks), tensor calculus, backpropagation gradients, and gating dynamics are mathematically formalized for Time-Series Forecasting Using LSTM Networks. \\
    \addlinespace
\textbf{PO2} & \textbf{Problem Analysis:} Systematically identifies, formulates, and investigates complex computational challenges, diagnosing architectural bottlenecks, concurrency anomalies, latency constraints, and integrity failures through rigorous analytical techniques. In this case study (Time-Series Forecasting Using LSTM Networks), systematic problem analysis resolves vanishing gradients, spatial occlusion, and high-frequency noise in deep representations. \\
    \addlinespace
\textbf{PO4} & \textbf{Conduct Investigations of Complex Problems:} Uses research-based knowledge and experimental methodologies, including synthesis of empirical benchmarks, stress testing, and profiling, to draw valid technical conclusions. In this case study (Time-Series Forecasting Using LSTM Networks), rigorous empirical validation using 5-fold cross-validation, ablation studies, and Grad-CAM interpretability provides experimental verification. \\
    \addlinespace
\textbf{PO5} & \textbf{Modern Tool Usage:} Selects and applies state-of-the-art software frameworks, modeling tools, database management engines, profiling utilities, and development environments to deliver high-performance solutions. In this case study (Time-Series Forecasting Using LSTM Networks), PyTorch, TensorRT, CUDA acceleration, and high-performance computing environments are proficiently utilized. \\
    \addlinespace
\textbf{PSO1} & \textbf{Program-Specific (AI/ML Core):} Leverages advanced intelligent computational techniques, algorithmic optimization, predictive analytics, and automated decision-making workflows to enhance system capabilities. In this case study (Time-Series Forecasting Using LSTM Networks), deep convolutional, recurrent, and transformer representations form the core of modern intelligent software systems. \\
    \addlinespace
\textbf{PSO2} & \textbf{Program-Specific (Applied Engineering):} Designs and implements production-grade end-to-end software pipelines integrating intelligent modules with robust backend infrastructure and enterprise services. In this case study (Time-Series Forecasting Using LSTM Networks), deployment of quantized neural networks onto embedded edge accelerators demonstrates end-to-end applied engineering excellence. \\
    \addlinespace
    \bottomrule
\end{longtable}}